\chapter{2009年大纲I卷（理）}

\section{单选题}

\begin{problem}
设集合 \(A = \{4, 5, 7, 9\}\)，\(B = \{3, 4, 7, 8, 9\}\)，全集 \(U = A \cup B\)，则集合 \(\complement_U\left(A \cap B\right)\) 中的元素共有（\quad）

\choices
    {\(3\text{ 个}\)}
    {\(4\text{ 个}\)}
    {\(5\text{ 个}\)}
    {\(6\text{ 个}\)}
\end{problem}

\begin{answer}
A
\end{answer}

\begin{solution}
由集合定义，
\(A\cup B=\{3,4,5,7,8,9\}\)，\(A\cap B=\{4,7,9\}\)，所以
\(\complement_U\left(A\cap B\right)=\{3,5,8\}\)，共有 \(3\) 个元素，故选 A.
\end{solution}

\begin{problem}
已知 \(\frac{\overline{z}}{1 + \i} = 2 + \i\)，则复数 \(z =\)（\quad）

\choices
    {\(-1 + 3\i\)}
    {\(1 - 3\i\)}
    {\( 3 + \i \)}
    {\(3 - \i\)}
\end{problem}

\begin{answer}
B
\end{answer}

\begin{solution}
由题设
\(\overline{z}=(2+\i)(1+\i)=1+3\i\)，两边取共轭得
\(z=1-3\i\)，故选 B.
\end{solution}

\begin{problem}
不等式 \(\left|\frac{x + 1}{x - 1}\right| < 1\) 的解集为（\quad）

\choices
    {\(\{x\mid 0 < x < 1\} \cup \{x\mid x > 1\}\)}
    {\(\{x\mid 0 < x < 1\}\)}
    {\(\{x\mid -1 < x < 0\}\)}
    {\(\{x\mid x < 0\}\)}
\end{problem}

\begin{answer}
D
\end{answer}

\begin{solution}
定义域要求 \(x\ne1\). 在此条件下，原不等式等价于
\(\lvert x+1\rvert<\lvert x-1\rvert\)，两边平方得
\(x^2+2x+1<x^2-2x+1\)，即 \(x<0\). 因此解集为
\(\{x\mid x<0\}\)，故选 D.
\end{solution}

\begin{problem}
设双曲线 \(\frac{x^2}{a^2} - \frac{y^2}{b^2} = 1 \) （\(a > 0\)，\(b > 0\)）的渐近线与抛物线 \(y = x^2 + 1\) 相切，则该双曲线的离心率等于（\quad）

\choices
    {\(\sqrt{3}\)}
    {\(2\)}
    {\(\sqrt{5}\)}
    {\(\sqrt{6}\)}
\end{problem}

\begin{answer}
C
\end{answer}

\begin{solution}
双曲线的渐近线为 \(y=\pm\frac{b}{a}x\). 令 \(k=\frac{b}{a}>0\)，将其中一条渐近线代入
\(y=x^2+1\)，得到二次方程 \(x^2\mp kx+1=0\). 相切时判别式为零，所以
\(k^2-4=0\)，从而 \(\frac{b}{a}=2\). 双曲线的焦半距满足
\(c^2=a^2+b^2\)，故离心率
\(e=\frac{c}{a}=\sqrt{1+\frac{b^2}{a^2}}=\sqrt5\)，故选 C.
\end{solution}

\begin{problem}
甲组有 \(5\text{ 名}\)男同学，\(3\text{ 名}\)女同学；乙组有 \(6\text{ 名}\)男同学，\(2\text{ 名}\)女同学. 若从甲，乙两组中各选出 \(2\text{ 名}\)同学，则选出的 \(4\text{ 人}\)中恰有 \(1\text{ 名}\)女同学的不同选法共有（\quad）

\choices
    {\(150\text{ 种}\)}
    {\(180\text{ 种}\)}
    {\(300\text{ 种}\)}
    {\(345\text{ 种}\)}
\end{problem}

\begin{answer}
D
\end{answer}

\begin{solution}
恰有一名女生分两类计数. 若女生来自甲组，则选法数为
\(\mathrm C_3^1\mathrm C_5^1\mathrm C_6^2=3\times5\times15=225\)；若女生来自乙组，则选法数为
\(\mathrm C_5^2\mathrm C_2^1\mathrm C_6^1=10\times2\times6=120\). 两类互斥，所以总数为
\(225+120=345\)，故选 D.
\end{solution}

\begin{problem}
设 \(\bs{a}\)，\(\bs{b}\)，\(\bs{c}\) 是单位向量，且 \(\bs{a} \cdot \bs{b} = 0\)，则 \(\left(\bs{a} - \bs{c}\right) \cdot \left(\bs{b} - \bs{c}\right)\) 的最小值为（\quad）

\choices
    {\(-2\)}
    {\(\sqrt{2} - 2\)}
    {\(-1\)}
    {\(1 - \sqrt{2}\)}
\end{problem}

\begin{answer}
D
\end{answer}

\begin{solution}
利用单位向量条件展开原式：
\[
(\bs{a}-\bs{c})\cdot(\bs{b}-\bs{c})=\bs{a}\cdot\bs{b}-\bs{a}\cdot\bs{c}-\bs{b}\cdot\bs{c}+\bs{c}\cdot\bs{c}=1-\bs{c}\cdot(\bs{a}+\bs{b})
\]
又因为 \(\bs{a}\perp\bs{b}\) 且二者均为单位向量，
\(\lvert\bs{a}+\bs{b}\rvert=\sqrt2\). 由数量积不等式，
\(\bs{c}\cdot(\bs{a}+\bs{b})\leqslant\lvert\bs{c}\rvert\lvert\bs{a}+\bs{b}\rvert=\sqrt2\)，因此原式不小于 \(1-\sqrt2\). 当 \(\bs{c}=\frac{\bs{a}+\bs{b}}{\sqrt2}\) 时取等号，所以最小值为 \(1-\sqrt2\)，故选 D.
\end{solution}

\begin{problem}
已知三棱柱 \(ABC - A_{1}B_{1}C_{1}\) 的侧棱与底面边长都相等，\(A_{1}\) 在底面 \(ABC\) 上的射影为 \(BC\) 的中点，则异面直线 \(AB\) 与 \(CC_{1}\) 所成的角的余弦值为（\quad）

\choices
    {\(\frac{\sqrt{3}}{4}\)}
    {\(\frac{\sqrt{5}}{4}\)}
    {\(\frac{\sqrt{7}}{4}\)}
    {\(\frac{3}{4}\)}
\end{problem}

\begin{answer}
D
\end{answer}

\begin{solution}
设底面边长为 \(s>0\)，令 \(D\) 为 \(BC\) 的中点，取坐标
\(B\left(-\frac{s}{2},0,0\right)\)，\(C\left(\frac{s}{2},0,0\right)\)，
\(A\left(0,\frac{\sqrt3}{2}s,0\right)\)，\(D=(0,0,0)\). 由于 \(A_1\) 在底面上的射影为 \(D\)，可设
\(A_1=\left(0,0,h\right)\). 由 \(AA_1=s\) 和
\(AD=\frac{\sqrt3}{2}s\)，得
\(h^2=s^2-\frac34s^2=\frac14s^2\)，故可取 \(h=\frac{s}{2}\).

三棱柱的侧棱向量为
\(\overrightarrow{CC_1}=\overrightarrow{AA_1}=\left(0,-\frac{\sqrt3}{2}s,\frac{s}{2}\right)\)，而
\(\overrightarrow{AB}=\left(-\frac{s}{2},-\frac{\sqrt3}{2}s,0\right)\). 因而
\(\overrightarrow{AB}\cdot\overrightarrow{CC_1}=\frac34s^2\)，\(\lvert\overrightarrow{AB}\rvert=\lvert\overrightarrow{CC_1}\rvert=s\).
异面直线所成角取锐角，故其余弦为 \(\frac34\)，故选 D.
\end{solution}

\begin{problem}
如果函数 \( y = 3\cos\left(2x + \varphi\right) \) 的图象关于点 \( \left(\frac{4\pi}{3}, 0\right) \) 中心对称，那么 \( |\varphi| \) 的最小值为（\quad）

\choices
    {\( \frac{\pi}{6} \)}
    {\( \frac{\pi}{4} \)}
    {\( \frac{\pi}{3} \)}
    {\( \frac{\pi}{2} \)}
\end{problem}

\begin{answer}
A
\end{answer}

\begin{solution}
周期函数 \(y=3\cos(2x+\varphi)\) 的图象在函数值为零的点处具有中心对称性. 因而令
\(x_0=\frac{4\pi}{3}\)，有
\(2x_0+\varphi=\frac{\pi}{2}+k\pi\)，\((k\in\mathbb Z)\).
确实，当该条件成立时，对任意 \(t\) 有
\(3\cos(2(x_0+t)+\varphi)=-3\cos(2(x_0-t)+\varphi)\)，所以中心为
\((x_0,0)\).
所以
\(\varphi=-\frac{13\pi}{6}+k\pi\). 取 \(k=2\) 得 \(\varphi=-\frac{\pi}{6}\)，且所有可行值相差 \(\pi\)，故 \(\lvert\varphi\rvert\) 的最小值为 \(\frac{\pi}{6}\)，故选 A.
\end{solution}

\begin{problem}
已知直线 \( y = x + 1 \) 与曲线 \( y = \ln\left(x + a\right) \) 相切，则 \( a \) 的值为（\quad）

\choices
    {\(1\)}
    {\(2\)}
    {\(-1\)}
    {\(-2\)}
\end{problem}

\begin{answer}
B
\end{answer}

\begin{solution}
设切点为 \(P(x_0,y_0)\). 曲线 \(y=\ln(x+a)\) 在切点处的导数为 \(\frac{1}{x_0+a}\)，而切线 \(y=x+1\) 的斜率为 \(1\)，故
\(x_0+a=1\). 于是 \(y_0=\ln1=0\)，又因 \(P\) 在直线上，\(x_0+1=0\)，所以 \(x_0=-1\)，从而 \(a=2\)，故选 B.
\end{solution}

\begin{problem}
已知二面角 \(\alpha - l - \beta\) 为 \(60^{\circ}\)，动点 \(P\)，\(Q\) 分别在面 \(\alpha\)，\(\beta\) 内，\(P\) 到 \(\beta\) 的距离为 \(\sqrt{3}\)，\(Q\) 到 \(\alpha\) 的距离为 \(2\sqrt{3}\)，则 \(P\)，\(Q\) 两点之间距离的最小值为（\quad）

\choices
    {\(\sqrt{2}\)}
    {\(2\)}
    {\(2\sqrt{3}\)}
    {\(4\)}
\end{problem}

\begin{answer}
C
\end{answer}

\begin{solution}
在垂直于棱 \(l\) 的截面内，面 \(\alpha\)、\(\beta\) 截成夹角为 \(60^{\circ}\) 的两条射线. 设 \(P\)、\(Q\) 到棱 \(l\) 的距离分别为 \(s\)、\(t\). 由点到对面距离，
\(s\sin60^{\circ}=\sqrt3\)，\(t\sin60^{\circ}=2\sqrt3\)，所以 \(s=2\)，\(t=4\).

若 \(P\)、\(Q\) 在棱方向上的坐标差为 \(h\)，则由余弦定理
\[
PQ^2=h^2+s^2+t^2-2st\cos60^{\circ}=h^2+12\geqslant12
\]
当二者在同一垂直于 \(l\) 的截面内时 \(h=0\)，等号成立，因此 \(PQ\) 的最小值为 \(2\sqrt3\)，故选 C.
\end{solution}

\begin{problem}
函数 \( f(x) \) 的定义域为 \( \R \)，若 \( f(x + 1) \) 与 \( f(x - 1) \) 都是奇函数，则（\quad）

\choices
    {\(f(x)\) 是偶函数}
    {\(f(x)\) 是奇函数}
    {\(f(x) = f(x + 2)\)}
    {\(f(x + 3)\) 是奇函数}
\end{problem}

\begin{answer}
D
\end{answer}

\begin{solution}
由 \(f(x+1)\) 为奇函数，令 \(t=x+1\)，得
\(f(t)=-f(2-t)\). 由 \(f(x-1)\) 为奇函数，令 \(t=x-1\)，得
\(f(t)=-f(-2-t)\). 两式比较可得 \(f(2-t)=f(-2-t)\)，即
\(f(u)=f(u-4)\)，所以 \(4\) 是 \(f\) 的一个周期.

再由 \(f(t)=-f(2-t)\)，取 \(t=3-x\)，有
\(f(3-x)=-f(x-1)=-f(x+3)\)，其中最后一个等号利用了周期 \(4\). 因而
\(f(x+3)\) 为奇函数，故选 D.
\end{solution}

\begin{problem}
已知椭圆 \(C: \frac{x^2}{2} + y^2 = 1\) 的右焦点为 \(F\)，右准线为 \(l\)，点 \(A \in l\)，线段 \(AF\) 交 \(C\) 于点 \(B\). 若 \(\overrightarrow{FA} = 3\overrightarrow{FB}\)，则 \(\left|AF\right| =\)（\quad）

\choices
    {\(\sqrt{2}\)}
    {\(2\)}
    {\(\sqrt{3}\)}
    {\(3\)}
\end{problem}

\begin{answer}
A
\end{answer}

\begin{solution}
椭圆的长半轴为 \(\sqrt2\)，短半轴为 \(1\)，故右焦点为 \(F=(1,0)\)，右准线为 \(x=2\). 设 \(A=(2,t)\). 由
\(\overrightarrow{FA}=3\overrightarrow{FB}\)，得
\(B=F+\frac13(A-F)=\left(\frac43,\frac{t}{3}\right)\).
将 \(B\) 代入椭圆方程，得
\[
\frac{(4/3)^2}{2}+\left(\frac{t}{3}\right)^2=1
\]
即 \(t^2=1\). 因此
\(\lvert AF\rvert^2=(2-1)^2+t^2=2\)，所以 \(\lvert AF\rvert=\sqrt2\)，故选 A.
\end{solution}

\section{填空题}

\begin{problem}
\((x - y)^{10}\) 的展开式中，\(x^{7}y^{3}\) 的系数与 \(x^{3}y^{7}\) 的系数之和等于\fillinblank{}.
\end{problem}

\begin{answer}
\(-240\)
\end{answer}

\begin{solution}
二项式通项为 \(\mathrm C_{10}^{k}x^{10-k}(-y)^k\). 其中 \(x^7y^3\) 对应 \(k=3\)，系数为
\(-\mathrm C_{10}^{3}=-120\)；\(x^3y^7\) 对应 \(k=7\)，系数为
\(-\mathrm C_{10}^{7}=-120\). 所以两系数之和为 \(-120-120=-240\).
\end{solution}

\begin{problem}
设等差数列 \(\{a_{n}\}\) 的前 \(n\) 项和为 \(S_{n}\)，若 \(S_{9} = 72\)，则 \(a_{2} + a_{4} + a_{9} = \)\fillinblank{}.
\end{problem}

\begin{answer}
\(24\)
\end{answer}

\begin{solution}
设等差数列首项为 \(a_1\)，公差为 \(d\). 由
\[
S_9=\frac{9}{2}(a_1+a_9)=\frac{9}{2}(2a_1+8d)=9(a_1+4d)=72
\]
得 \(a_1+4d=8\)，即 \(a_5=8\). 又
\[
a_2+a_4+a_9=(a_1+d)+(a_1+3d)+(a_1+8d)=3(a_1+4d)=3a_5=24
\]
\end{solution}

\begin{problem}
直三棱柱 \( ABC - A_{1}B_{1}C_{1} \) 的各顶点都在同一球面上. 若 \(AB = AC = AA_{1} = 2\)，\(\angle BAC = 120^{\circ}\)，则此球的表面积等于\fillinblank{}.
\end{problem}

\begin{answer}
\(20\pi\)
\end{answer}

\begin{solution}
由余弦定理，底面三角形满足
\[
BC^2=AB^2+AC^2-2AB\cdot AC\cos120^{\circ}=4+4+4=12
\]
所以 \(BC=2\sqrt3\). 由正弦定理，底面三角形外接圆半径为
\(R_0=\frac{BC}{2\sin120^{\circ}}=2\).

直三棱柱的球心在上下两个底面外接圆圆心连线的中点处，球半径 \(R\) 满足
\[
R^2=R_0^2+\left(\frac{AA_1}{2}\right)^2=2^2+1^2=5
\]
故球的表面积为 \(4\pi R^2=20\pi\).
\end{solution}

\begin{problem}
若 \(\frac{\pi}{4} < x < \frac{\pi}{2}\)，则函数 \(y = \tan 2x \tan^3 x\) 的最大值为\fillinblank{}.
\end{problem}

\begin{answer}
\(-8\)
\end{answer}

\begin{solution}
在 \(\frac{\pi}{4}<x<\frac{\pi}{2}\) 内令 \(t=\tan x\)，则 \(t>1\)，且
\[
y=\tan2x\tan^3x=\frac{2t}{1-t^2}t^3=-\frac{2t^4}{t^2-1}
\]
令 \(u=t^2-1>0\)，则
\[
y=-2\left(u+2+\frac1u\right)\leqslant-2(2+2)=-8
\]
其中使用了 \(u+\frac1u\geqslant2\). 当 \(u=1\)，即 \(t=\sqrt2\) 时取等号，且此时的 \(x\) 在给定区间内. 所以最大值为 \(-8\).
\end{solution}

\section{解答题}

\begin{problem}
在 \(\triangle ABC\) 中，内角 \(A\)，\(B\)，\(C\) 的对边长分别为 \(a\)，\(b\)，\(c\)，已知 \(a^2 - c^2 = 2b\)，且 \(\sin A \cos C = 3 \cos A \sin C\)，求 \(b\).
\end{problem}

\begin{solution}
由正弦定理和余弦定理，得 \(\sin A=\frac{a}{2R}\)，\(\sin C=\frac{c}{2R}\)，\(\cos C=\frac{a^2+b^2-c^2}{2ab}\)，\(\cos A=\frac{b^2+c^2-a^2}{2bc}\)，
其中 \(R\) 为三角形 \(ABC\) 的外接圆半径. 将它们代入题设等式，得
\[
\frac{a^2+b^2-c^2}{4bR}
=3\frac{b^2+c^2-a^2}{4bR}
\]
所以
\[
a^2+b^2-c^2=3\left(b^2+c^2-a^2\right)
\]
即 \(2a^2-b^2-2c^2=0\)，从而
\[
b^2=2\left(a^2-c^2\right)
\]
又由 \(a^2-c^2=2b\)，得 \(b^2=4b\). 因为 \(b\) 是三角形的边长，\(b>0\)，所以
\[
b=4
\]
\end{solution}

\begin{problem}
如图，四棱锥 \(S - ABCD\) 中，底面 \(ABCD\) 为矩形，\(SD \perp\) 底面 \(ABCD\)，\(AD = \sqrt{2}\)，\(DC = SD = 2\)，点 \(M\) 在侧棱 \(SC\) 上，\(\angle ABM = 60^\circ\).

\begin{enumerate}
\item 证明：\(M\) 是侧棱 \(SC\) 的中点；
\item 求二面角 \(S - AM - B\) 的大小.
\end{enumerate}

\bitmapfigure[max width=0.42\linewidth]{img/2009/syllabus_paper_1_science/q18_fig1.png}
\end{problem}

\begin{solution}
建立空间直角坐标系，令 \(D(0,0,0)\)，\(C(2,0,0)\)，\(A(0,\sqrt{2},0)\)，\(B(2,\sqrt{2},0)\)，\(S(0,0,2)\).
设 \(M\) 将 \(SC\) 按参数 \(\lambda\) 分割，则 \(M\left(2\lambda,0,2-2\lambda\right)\)，其中 \(0\leqslant\lambda\leqslant1\).
于是 \(\overrightarrow{BA}=(-2,0,0)\)，\(\overrightarrow{BM}=\left(2\lambda-2,-\sqrt{2},2-2\lambda\right)\)，且 \(\overrightarrow{BA}\cdot\overrightarrow{BM}=4(1-\lambda)\)，\(\left|\overrightarrow{BM}\right|^2=8(1-\lambda)^2+2\).
由 \(\angle ABM=60^\circ\)，得
\[
\frac{4(1-\lambda)}{2\sqrt{8(1-\lambda)^2+2}}=\cos60^\circ=\frac{1}{2}
\]
由于 \(1-\lambda\geqslant0\)，平方并整理得
\[
16(1-\lambda)^2=8(1-\lambda)^2+2
\]
从而 \((1-\lambda)^2=\frac{1}{4}\)，结合 \(1-\lambda\geqslant0\) 得 \(\lambda=\frac{1}{2}\).
所以 \(M(1,0,1)\)，即 \(M\) 是侧棱 \(SC\) 的中点.

设二面角 \(S-AM-B\) 的大小为 \(\theta\). 取公共棱 \(AM\) 的方向向量 \(\bs{u}=\overrightarrow{AM}=(1,-\sqrt{2},1)\)，则 \(\left|\bs{u}\right|^2=4\).
在平面 \(SAM\) 和平面 \(BAM\) 内，分别取指向 \(S\) 和 \(B\) 一侧且垂直于 \(AM\) 的向量
\(\bs{p}=\overrightarrow{AS}-\bs{u}=(-1,0,1)\)，
\(\bs{q}=\overrightarrow{AB}-\frac{1}{2}\bs{u}=\left(\frac{3}{2},\frac{\sqrt{2}}{2},-\frac{1}{2}\right)\).
其中 \(\overrightarrow{AS}=(0,-\sqrt{2},2)\)，\(\overrightarrow{AB}=(2,0,0)\)，且 \(\frac{\overrightarrow{AS}\cdot\bs{u}}{\left|\bs{u}\right|^2}=1\)，\(\frac{\overrightarrow{AB}\cdot\bs{u}}{\left|\bs{u}\right|^2}=\frac{1}{2}\)，所以 \(\bs{p}\) 与 \(\bs{q}\) 均垂直于 \(\bs{u}\). 因而 \(\bs{p}\cdot\bs{q}=-2\)，\(\left|\bs{p}\right|=\sqrt{2}\)，\(\left|\bs{q}\right|=\sqrt{3}\).
所以
\[
\cos\theta=\frac{\bs{p}\cdot\bs{q}}{\left|\bs{p}\right|\left|\bs{q}\right|}
=-\frac{2}{\sqrt{2}\sqrt{3}}=-\frac{\sqrt{6}}{3}
\]
故二面角 \(S-AM-B\) 的大小为
\[
\theta=\arccos\left(-\frac{\sqrt{6}}{3}\right)
\]
\end{solution}

\begin{problem}
甲，乙二人进行一次围棋比赛，约定先胜 \(3\text{ 局}\)者获得这次比赛的胜利，比赛结束. 假设在一局中，甲获胜的概率为 \(0.6\)，乙获胜的概率为 \(0.4\)，各局比赛结果相互独立，已知前 \(2\text{ 局}\)中，甲，乙各胜 \(1\text{ 局}\).
\begin{enumerate}
\item 求甲获得这次比赛胜利的概率；
\item 设 \(\xi\) 表示从第 \(3\text{ 局}\)开始到比赛结束所进行的局数，求 \(\xi\) 的分布列及数学期望.
\end{enumerate}
\end{problem}

\begin{solution}
前 \(2\) 局中甲、乙各胜 \(1\) 局，所以从第 \(3\) 局开始，甲、乙都还需要胜 \(2\) 局才能获胜. 记甲、乙在一局中获胜的概率分别为 \(p=0.6=\frac{3}{5}\)，\(q=0.4=\frac{2}{5}\).

\begin{enumerate}
\item 甲获胜的互斥情形为：第 \(3\)、\(4\) 局连胜；第 \(3\) 局乙胜且第 \(4\)、\(5\) 局甲连胜；第 \(3\) 局甲胜、第 \(4\) 局乙胜且第 \(5\) 局甲胜. 因此
\[
P(\text{甲获胜})=p^2+2p\cdot q\cdot p=p^2+2p^2q
=\left(\frac{3}{5}\right)^2+2\left(\frac{3}{5}\right)^2\frac{2}{5}
=\frac{81}{125}
\]

\item 从第 \(3\) 局开始的前 \(2\) 局若均由同一方获胜，比赛在第 \(4\) 局结束；若这 \(2\) 局甲、乙各胜 \(1\) 局，则还需进行第 \(5\) 局. 因此
\[
P(\xi=2)=p^2+q^2
=\left(\frac{3}{5}\right)^2+\left(\frac{2}{5}\right)^2
=\frac{13}{25}
\]
\[
P(\xi=3)=2pq
=2\cdot\frac{3}{5}\cdot\frac{2}{5}
=\frac{12}{25}
\]
除 \(2\)，\(3\) 外，\(\xi\) 不取其他值，即
\[
P(\xi=k)=
\begin{cases}
\dfrac{13}{25},&k=2,\\
\dfrac{12}{25},&k=3,\\
0,&\text{其他}
\end{cases}
\]
所以
\[
E(\xi)=2\cdot\frac{13}{25}+3\cdot\frac{12}{25}=\frac{62}{25}
\]
\end{enumerate}
\end{solution}

\begin{problem}
在数列 \(\{a_{n}\}\) 中，\(a_{1} = 1\)，\(a_{n+1} = \left(1 + \frac{1}{n}\right)a_{n} + \frac{n+1}{2^{n}}\).
\begin{enumerate}
\item 设 \(b_{n}=\frac{a_{n}}{n}\)，求数列 \(\{b_{n}\}\) 的通项公式；
\item 求数列 \(\{a_{n}\}\) 的前 \(n\) 项和 \(S_{n}\).
\end{enumerate}
\end{problem}

\begin{solution}
\begin{enumerate}
\item 由 \(b_n=\frac{a_n}{n}\)，得 \(a_n=nb_n\)，\(a_{n+1}=(n+1)b_{n+1}\). 将它们代入递推式，得
\[
(n+1)b_{n+1}=\left(1+\frac{1}{n}\right)nb_n+\frac{n+1}{2^n}
=(n+1)b_n+\frac{n+1}{2^n}
\]
两边除以 \(n+1\)，得
\[
b_{n+1}=b_n+\frac{1}{2^n}
\]
又 \(b_1=\frac{a_1}{1}=1\)，所以对 \(n\geqslant2\)，
\[
b_n=b_1+\sum_{k=1}^{n-1}\frac{1}{2^k}
=1+\left(1-\frac{1}{2^{n-1}}\right)
=2-\frac{1}{2^{n-1}}
\]
该式在 \(n=1\) 时也成立，故
\[
b_n=2-\frac{1}{2^{n-1}}
\]

\item 由上式，得
\[
a_n=nb_n=2n-\frac{n}{2^{n-1}}
\]
因此
\[
S_n=\sum_{k=1}^{n}a_k
=n(n+1)-\sum_{k=1}^{n}\frac{k}{2^{k-1}}
\]
令
\[
T_n=\sum_{k=1}^{n}\frac{k}{2^{k-1}}
\]
则
\[
T_n=1+\frac{2}{2}+\frac{3}{2^2}+\cdots+\frac{n}{2^{n-1}}
\]
\[
\frac{1}{2}T_n=\frac{1}{2}+\frac{2}{2^2}+\frac{3}{2^3}+\cdots+\frac{n}{2^n}
\]
两式相减，得
\[
\frac{1}{2}T_n=1+\frac{1}{2}+\frac{1}{2^2}+\cdots+\frac{1}{2^{n-1}}-\frac{n}{2^n}
=2-\frac{n+2}{2^n}
\]
所以
\[
T_n=4-\frac{n+2}{2^{n-1}}
\]
从而
\[
S_n=n(n+1)-4+\frac{n+2}{2^{n-1}}
\]
\end{enumerate}
\end{solution}

\begin{problem}
如图，已知抛物线 \(E: y^{2} = x\) 与圆 \(M: \left(x - 4\right)^{2} + y^{2} = r^{2}\) （\(r > 0\)）相交于 \(A\)，\(B\)，\(C\)，\(D\) 四个点.

\begin{enumerate}
\item 求 \(r\) 的取值范围；
\item 当四边形 \(ABCD\) 的面积最大时，求对角线 \(AC\)，\(BD\) 的交点 \(P\) 坐标.
\end{enumerate}

\FigureLayoutDeclare{img/2009/syllabus_paper_1_science/q21_fig1.png}{10.0cm}{49bp 21bp 32bp 13bp}
\bitmapfigure[max width=0.42\linewidth]{img/2009/syllabus_paper_1_science/q21_fig1.png}
\end{problem}

\begin{solution}
\begin{enumerate}
\item 将抛物线 \(y^2=x\) 代入圆的方程，得
\[
\left(x-4\right)^2+x=r^2
\]
即
\[
x^2-7x+16-r^2=0
\]
要使圆与抛物线有四个交点，上式必须有两个不相等的正根. 其判别式为
\[
\Delta=49-4\left(16-r^2\right)=4r^2-15>0
\]
且两根之积为 \(16-r^2>0\). 因为两根之和为 \(7>0\)，两根均为正数，所以 \(r>\frac{\sqrt{15}}{2}\)，\(r<4\)，因此 \(r\) 的取值范围为 \(\frac{\sqrt{15}}{2}<r<4\).

\item 设上述二次方程的两个根为 \(x_1\)，\(x_2\)，且 \(0<x_1<x_2\). 设 \(u=\sqrt{x_1}\)，\(v=\sqrt{x_2}\)，则四个交点可表示为 \(A(x_1,u)\)，\(B(x_1,-u)\)，\(D(x_2,v)\)，\(C(x_2,-v)\).
由韦达定理，得 \(x_1+x_2=7\)，所以 \(u^2+v^2=7\). 四边形 \(ABCD\) 的两条平行边为 \(AB=2u\)，\(CD=2v\)，两平行边之间的距离为 \(x_2-x_1=v^2-u^2\)，故其面积为
\[
S=\frac{1}{2}\left(2u+2v\right)\left(v^2-u^2\right)=\left(u+v\right)^2(v-u)
\]
令 \(d=v-u\)，则 \(d>0\)，并且
\[
\left(u+v\right)^2+\left(v-u\right)^2=2\left(u^2+v^2\right)=14
\]
于是
\[
S=\left(14-d^2\right)d=14d-d^3
\]
当 \(0<d<\sqrt{7}\) 时，
\[
S'=14-3d^2
\]
所以当 \(d^2=\frac{14}{3}\) 时，\(S\) 取得最大值. 此时
\[
\left(u+v\right)^2=14-\frac{14}{3}=\frac{28}{3}
\]

由图形关于 \(x\) 轴对称，且两条对角线关于 \(x\) 轴互为对称图形，故其交点 \(P\) 在 \(x\) 轴上. 直线 \(AC\) 与 \(x\) 轴交点的横坐标为
\[
x_P=x_1+\frac{x_2-x_1}{u+v}u
=u^2+\frac{(v-u)(v+u)}{u+v}u
=uv
\]
又
\[
uv=\frac{(u+v)^2-(v-u)^2}{4}
=\frac{\frac{28}{3}-\frac{14}{3}}{4}=\frac{7}{6}
\]
因此
\[
P\left(\frac{7}{6},0\right)
\]
\end{enumerate}
\end{solution}

\begin{problem}
设函数 \( f(x) = x^3 + 3bx^2 + 3cx \) 有两个极值点 \(x_1\)，\(x_2\)，且 \( x_1 \in [-1, 0] \)，\( x_2 \in [1, 2] \).

\begin{enumerate}
\item 求 \(b\)，\(c\) 满足的约束条件，并在下面的坐标平面内，画出满足这些条件的点 \((b, c)\) 的区域；
\item 证明：\(-10 \leqslant f(x_{2}) \leqslant -\frac{1}{2}\).
\end{enumerate}

\bitmapfigure[max width=0.42\linewidth]{img/2009/syllabus_paper_1_science/q22_fig1.png}
\end{problem}

\begin{solution}
\begin{enumerate}
\item 由
\[
f'(x)=3x^2+6bx+3c=3\left(x^2+2bx+c\right)
\]
且 \(x_1\)，\(x_2\) 是两个极值点，知 \(f'(x)=3\left(x-x_1\right)\left(x-x_2\right)\)，其中 \(x_1\in[-1,0]\)，\(x_2\in[1,2]\).
因为二次函数 \(f'(x)\) 的开口向上，所以 \(f'(-1)\geqslant0\)，\(f'(0)\leqslant0\)，\(f'(1)\leqslant0\)，\(f'(2)\geqslant0\). 将 \(f'(x)=3x^2+6bx+3c\) 代入，得到 \(c-2b+1\geqslant0\)，\(c\leqslant0\)，\(c+2b+1\leqslant0\)，\(c+4b+4\geqslant0\).
即 \(b\)，\(c\) 满足的约束条件为 \(c\geqslant2b-1\)，\(c\leqslant0\)，\(c\leqslant-2b-1\)，\(c\geqslant-4b-4\).
反之，若上述四个不等式成立，则因 \(f'(x)\) 为开口向上的二次函数，且 \(f'(-1)\geqslant0\geqslant f'(0)\)、\(f'(1)\leqslant0\leqslant f'(2)\)，方程 \(f'(x)=0\) 在区间 \([-1,0]\)、\([1,2]\) 内各有一个根；两区间不相交，故这两个根正是两个极值点. 因此，满足条件的点 \((b,c)\) 构成由直线 \(c=0\)、\(c=-2b-1\)、\(c=2b-1\)、\(c=-4b-4\) 围成的四边形区域，其四个顶点依次为 \((-1,0)\)、\(\left(-\frac{1}{2},0\right)\)、\((0,-1)\)、\(\left(-\frac{1}{2},-2\right)\).

\item 由韦达定理，得 \(x_1+x_2=-2b\)，\(x_1x_2=c\)，所以 \(b=-\frac{x_1+x_2}{2}\)，\(c=x_1x_2\).
在 \(x=x_2\) 处，
\[
f(x_2)=x_2^3+3bx_2^2+3cx_2
=x_2^3-\frac{3}{2}(x_1+x_2)x_2^2+3x_1x_2^2
=\frac{x_2^2}{2}\left(3x_1-x_2\right)
\]
因为 \(x_1\in[-1,0]\)，\(x_2\in[1,2]\)，且 \(x_2^2/2>0\)，所以
\[
-\frac{x_2^2}{2}\left(x_2+3\right)
\leqslant f(x_2)
\leqslant -\frac{x_2^3}{2}
\]
对左端，函数 \(-\frac{1}{2}x_2^2\left(x_2+3\right)\) 在 \([1,2]\) 上递减，因此其最小值为
\[
-\frac{1}{2}\cdot2^2\left(2+3\right)=-10
\]
对右端，函数 \(-\frac{x_2^3}{2}\) 在 \([1,2]\) 上递减，因此其最大值为
\[
-\frac{1^3}{2}=-\frac{1}{2}
\]
于是
\[
-10\leqslant f(x_2)\leqslant-\frac{1}{2}
\]
\end{enumerate}
\end{solution}
